% Standalone reference document. Build from the repository root with make examples.
% Shared document class, packages, and formatting for LuaLaTeX.
% Based on the TUDoThesis template by Maximilian Nöthe; see LICENSE.
\documentclass[
  tucolor,               % Remove for headings without TU green.
  BCOR=12mm,             % Adjust the binding correction for the final binding.
  parskip=half,
  open=any,
  cleardoublepage=plain,
]{tudothesis}

\setcounter{tocdepth}{1} % List chapters and sections in the contents.
\usepackage[aux]{rerunfilecheck}

% Fonts and languages. The last Babel language is the main language.
\usepackage{fontspec}
\defaultfontfeatures{Ligatures=TeX}
\usepackage[ngerman, american]{babel}
\usepackage[autostyle]{csquotes}
\usepackage{microtype}

% Mathematics. unicode-math supplies the symbols formerly loaded by amssymb.
\usepackage{mathtools}
\usepackage[
  math-style=ISO,
  bold-style=ISO,
  sans-style=italic,
  nabla=upright,
  partial=upright,
  warnings-off={mathtools-colon,mathtools-overbracket},
]{unicode-math}
\setmathfont{Latin Modern Math}
\usepackage{xfrac}

% Numbers and units follow the active document language.
\usepackage[
  locale=US,
  separate-uncertainty=true,
  per-mode=symbol-or-fraction,
]{siunitx}
\addto\extrasamerican{\sisetup{locale=US}}
\addto\extrasngerman{\sisetup{locale=DE}}

% Tables and graphics. Caption/subcaption support is provided by tudothesis.
\usepackage{booktabs}
\usepackage{graphicx}
\graphicspath{{figures/}{plots/}{logos/}}
\usepackage{scrhack}
\usepackage{float}
\floatplacement{figure}{htbp}
\floatplacement{table}{htbp}
\usepackage[section, below]{placeins}
\usepackage{pdfpages}
\pdfvariable minorversion=7

\usepackage{enumitem}
\setlist{nolistsep}
\usepackage[shortcuts]{extdash}

% Each document declares its own .bib resource after loading this header.
\usepackage[
  backend=biber,
  autolang=hyphen,
]{biblatex}

% Load hyperlink support near the end, before bookmarks and glossaries.
\usepackage[
  unicode,
  colorlinks=true,
  linkcolor=tugreen!55!black,
  citecolor=tugreen!55!black,
  urlcolor=blue!70!black,
]{hyperref}
\usepackage{bookmark}
\usepackage[toc, acronym]{glossaries}
\makenoidxglossaries % Uses LaTeX reruns; no external indexer is required.

% Reusable notation belongs here; package and layout settings belong in header.tex.
% Bold italic vectors, using unicode-math's ISO-style alphabet.
\newcommand{\vect}[1]{\symbf{#1}}

\addbibresource{examples/references.bib}
% Example definitions, loaded only by examples/formatting.tex.
\newacronym{ml}{ML}{machine learning}
\newglossaryentry{estimator}{
  name={estimator},
  description={a statistic used to estimate an unknown parameter},
}

\hypersetup{pdftitle={LaTeX Formatting Examples}, pdfauthor={}}

\begin{document}
\mainmatter
\chapter{Formatting Examples}
These examples share the thesis settings and are excluded from the thesis itself.
All numerical values below are illustrative.

\section{Text, citations, and acronyms}
Use \emph{emphasis}, \enquote{language-aware quotation marks}, and citations such as
\autocite{kuypers}. A citation in running text can use \textcite{ma-sebastian}.
Online resources can be cited in the same way \autocite{texmaker}.

The first use of \gls{ml} expands the acronym; later uses print \gls{ml}.
The term \gls{estimator} appears in the glossary.
Use \verb|\glspl{key}| for the plural of a defined term.

\section{Mathematics, numbers, and units}
\label{sec:example-math}
For a vector \(\vect{x} = (x_1, \ldots, x_N)\), the weighted mean is
\begin{equation}
  \bar{x} = \frac{\sum_{i=1}^{N} w_i x_i}{\sum_{i=1}^{N} w_i}.
  \label{eq:weighted-mean}
\end{equation}
Refer to it with Equation~\eqref{eq:weighted-mean}.
Inline fractions such as \(\sfrac{1}{2}\) are also supported.
Typeset measurements as \qty{1.23(4)}{\metre}, speeds as
\qty{3.0e8}{\metre\per\second}, and percentages as \qty{95}{\percent}.

\begin{otherlanguage}{ngerman}
Die Zahlenformatierung folgt der Sprache: \num{1234.56} und \qty{1.23(4)}{\metre}.
\end{otherlanguage}

Back in English, the same value is \num{1234.56}.

\section{Figures and references}
\begin{figure}
  \centering
  \includegraphics[width=0.45\linewidth]{logos/tu-logo.pdf}
  \caption{The template's TU Dortmund logo, used here to demonstrate image inclusion.}
  \label{fig:example-logo}
\end{figure}
Refer to Figure~\ref{fig:example-logo} and Section~\ref{sec:example-math}.
Place scientific figures in \texttt{figures/} and generated plots in \texttt{plots/}.
Use \verb|\includegraphics| with the filename; both directories are searched automatically.

\begin{figure}
  \centering
  \begin{subfigure}{0.45\linewidth}
    \centering
    \fbox{\rule{0pt}{15mm}\rule{0.8\linewidth}{0pt}}
    \caption{First panel.}
    \label{fig:example-panel-a}
  \end{subfigure}
  \hfill
  \begin{subfigure}{0.45\linewidth}
    \centering
    \fbox{\rule{0pt}{15mm}\rule{0.8\linewidth}{0pt}}
    \caption{Second panel.}
    \label{fig:example-panel-b}
  \end{subfigure}
  \caption{Replace the boxes with images to create a figure with two panels.}
  \label{fig:example-panels}
\end{figure}
Individual panels can be referenced as Figure~\ref{fig:example-panel-a}.

\section{Tables}
\begin{table}
  \centering
  \caption{Illustrative measurements. The \texttt{S} column aligns decimal markers.}
  \label{tab:example-measurements}
  \begin{tabular}{l S[table-format=2.2]}
    \toprule
    Sample & {Time / \unit{\second}} \\
    \midrule
    A & 1.23 \\
    B & 12.34 \\
    \bottomrule
  \end{tabular}
\end{table}
Table~\ref{tab:example-measurements} uses \texttt{booktabs} rules and places its caption above the data.

\backmatter
\printnoidxglossaries
\printbibliography
\end{document}
